%=====================================================================
\chapter{What Counts as a Contribution in Computing Research}\label{ch:contribution}
%=====================================================================
\locator{1}{Part~1 --- Foundations of Inquiry in Computing}

\begin{outcomes}
By the end of this chapter you will be able to:
\begin{tight}
  \item \textbf{Distinguish} research from development, and say which one a
    given piece of work is.
  \item \textbf{Classify} a computing paper into one of eight contribution
    types, and name the evidence each type owes its reader.
  \item \textbf{State} any proposed study as a claim--evidence--falsification
    triple.
  \item \textbf{Diagnose} the four ways a proposal fails to be a contribution
    at all.
  \item \textbf{Relate} your own topic to the relevance criteria your
    dissertation will be judged against.
\end{tight}
\end{outcomes}

\begin{prereq}
None. This is where the course starts. It assumes you have a degree in
computing and a vague idea of what you want to work on --- which is exactly the
situation this chapter is designed to fix.
\end{prereq}

\begin{keyterms}
contribution $\bullet$ claim $\bullet$ evidence $\bullet$ falsification
$\bullet$ artefact $\bullet$ empirical study $\bullet$ replication
$\bullet$ benchmark $\bullet$ delta $\bullet$ novelty $\bullet$ significance
$\bullet$ validation
\end{keyterms}

\section{The Sentence That Is Not Yet Research}

Almost every graduate scholar arrives with a sentence like one of these:

\begin{tight}
  \item ``I want to work on federated learning for healthcare.''
  \item ``My topic is blockchain in supply chain management.''
  \item ``I will apply deep learning to intrusion detection.''
\end{tight}

None of these is a research topic. Each is a \emph{subject area} --- a region
of the literature, not a question about it. You could spend three years inside
any of them and produce nothing a dissertation examiner would accept, because
none of them contains a claim that could turn out to be wrong.

The distance between that sentence and a dissertation is the distance this
chapter covers. It is shorter than it looks, and it is crossed by asking three
questions in order.

\begin{definitionbox}[Research contribution]
A \term{contribution} is a claim about the world, supported by evidence
gathered in a way that could have failed to support it, which somebody who
already knows the field did not know before.

\medskip
Every clause in that sentence is doing work. \emph{A claim}: not a topic, not
an activity, not a system. \emph{Supported by evidence}: not by plausibility or
enthusiasm. \emph{Could have failed}: if no possible result would have changed
your conclusion, you have not tested anything. \emph{Did not know before}: the
field, not you.
\end{definitionbox}

\subsection{Research is not development, and the difference is not difficulty}

Building a working federated-learning system for three hospitals is hard. It
may take a year, require real skill, and produce something genuinely useful.
It is still not, by itself, research --- and confusing the two is the single
most common reason a synopsis is sent back by a Departmental Research
Committee.

\begin{center}
\small
\begin{tabular}{@{}L{3.4cm}L{5.4cm}L{5.4cm}@{}}
\toprule
 & \textbf{Development} & \textbf{Research} \\
\midrule
Succeeds when      & The thing works
                   & The claim is tested, whichever way it comes out \\
A negative outcome & Is a failure
                   & Is a result, if the design was sound \\
Asks               & Can I build this?
                   & Is this true, and how would I know? \\
Ends with          & A deliverable
                   & A claim plus its evidence plus its limits \\
Is judged by       & Users
                   & Peers who know the alternatives \\
\bottomrule
\end{tabular}
\end{center}

The same project can be either. ``I built a federated-learning framework for
three hospitals'' is development. ``Federated training on non-IID clinical data
partitioned by hospital loses less accuracy under FedProx than under FedAvg,
and the gap widens with partition skew'' is a claim --- one that a careful
experiment could refute. The system is now the \emph{instrument}, not the
result.

\begin{conceptbox}[The instrument and the result]
Design science research (\chref{designscience}) complicates this usefully: there,
the artefact genuinely \emph{is} a contribution. But even then it contributes
by demonstrating something --- that a class of problem is solvable this way,
that a design principle transfers, that a trade-off can be relaxed. The
artefact still has to carry a claim. What is never a contribution is an
artefact accompanied only by the assertion that it was built.
\end{conceptbox}

\section{Eight Kinds of Contribution}

Computing research does not have one shape. A paper claiming a new algorithm
owes its reader completely different evidence from a paper claiming that
developers misunderstand static analysis warnings. Knowing which kind of thing
you are producing tells you what you must show.

\begin{center}
\small
\begin{longtable}{@{}L{3.0cm}L{5.2cm}L{6.0cm}@{}}
\toprule
\textbf{Type} & \textbf{The claim takes the form} & \textbf{Evidence it owes} \\
\midrule
\endfirsthead
\toprule
\textbf{Type} & \textbf{The claim takes the form} & \textbf{Evidence it owes} \\
\midrule
\endhead
\bottomrule
\endfoot
Technique or algorithm
  & ``This method achieves X under conditions C, better than the best prior
     method on dimension D''
  & Correctness argument or proof; comparison against a \emph{current} baseline
    on data the baseline was not tuned on; ablation showing which component
    earns the gain \\
Artefact / system
  & ``A system with properties P is achievable, and building it revealed
     design principle Q''
  & The artefact, plus an evaluation matched to the claim
    (\chref{designscience}); the principle stated separately from the code \\
Empirical study
  & ``In population P, X is associated with Y'' or ``treatment T changes
     outcome O''
  & A design that could have shown no effect; effect size with an interval,
    not only a p-value; threats to validity (\chref{validity}) \\
Theory or model
  & ``These phenomena are explained by mechanism M''
  & Derivation; predictions that differ from rival explanations; at least one
    confrontation with data \\
Methodology
  & ``Research of type R is better conducted this way''
  & Demonstration on real studies; comparison with existing practice; honest
    statement of where it does not apply \\
Dataset or benchmark
  & ``This resource enables questions that could not be asked before''
  & Construction and sampling documented; datasheet~\cite{gebru2021datasheets};
    baseline results; known biases and limits \\
Replication
  & ``The original finding does / does not hold under conditions C''
  & Faithfulness to the original protocol, or a documented and justified
    deviation; the same analysis, pre-specified \\
Survey / review
  & ``The state of knowledge on Q is S, and the gaps are G''
  & A protocol, a reproducible search, and a synthesis that is more than a
    list (\chref{slr}) \\
\end{longtable}
\end{center}

\begin{pitfall}[Owing the wrong evidence]
Most rejections come from delivering the evidence of one type while claiming
another.

\begin{tight}
  \item A \emph{system} paper that reports only accuracy numbers, with no
    design principle a reader could reuse --- the artefact was built, but
    nothing was learned.
  \item A \emph{technique} paper whose baseline is the 2016 method because the
    2024 one was harder to install. The comparison is not wrong, it is
    irrelevant.
  \item An \emph{empirical} paper reporting $p < 0.05$ on 12 participants with
    no effect size and no power analysis (\chref{sampling}).
  \item A \emph{survey} that tabulates 60 papers without ever answering a
    question. A table is not a synthesis.
\end{tight}
\end{pitfall}

\begin{paperbox}[Shaw on what reviewers actually look for]
Mary Shaw, \emph{Writing Good Software Engineering Research
Papers}, ICSE 2003~\cite{shaw2003writing}.

\medskip
\textbf{Why it matters.} Shaw analysed what ICSE accepted and rejected and
extracted the implicit rules: every paper answers a \emph{kind} of question,
produces a \emph{kind} of result, and validates it in a \emph{kind} of way ---
and reviewers reject papers where those three do not match each other. It
remains the clearest short statement of what a computing contribution is.

\medskip
\textbf{What to extract.} Her three tables of question types, result types and
validation types. Write down which cell of each your own project sits in. If
you cannot, your project is not yet defined.
\end{paperbox}

\section{The Claim--Evidence--Falsification Triple}

Here is the operational test. Before anything else --- before the literature
review, before choosing a method, certainly before writing code --- write three
lines:

\begin{center}
\small
\begin{tabular}{@{}L{2.6cm}L{11.7cm}@{}}
\toprule
\textbf{Claim} & One sentence, in the indicative, that could be true or false. \\
\textbf{Evidence} & What you would have to observe for a reasonable sceptic to
                    accept it. \\
\textbf{Falsifier} & What you could observe that would force you to withdraw
                     it. \\
\bottomrule
\end{tabular}
\end{center}

If the third line is empty, you do not have a research project. This is the
single most useful diagnostic in the course, and it costs five minutes.

\begin{worked}[From a subject area to a testable claim]
\textbf{Round 0.} ``I want to work on federated learning for healthcare.''
No claim. No falsifier. Not research.

\medskip
\textbf{Round 1.} ``I will build a federated learning system for hospitals in
Punjab.'' Now there is an activity, still no claim. The falsifier would be
``the system does not compile'', which nobody would publish.

\medskip
\textbf{Round 2.} ``Federated learning preserves privacy while maintaining
accuracy.'' A claim at last --- but not yours, and not testable as stated.
Preserves privacy against which adversary? Maintains accuracy relative to
what? This is a slogan from the abstract of somebody else's paper.

\medskip
\textbf{Round 3.} ``On clinical tabular data partitioned by hospital, FedProx
loses less accuracy than FedAvg as label skew increases.''
\begin{tight}
  \item \textbf{Claim:} stated, comparative, conditional on a data regime.
  \item \textbf{Evidence:} accuracy of both algorithms across a controlled
    range of skew, on more than one dataset, with repeated runs and intervals.
  \item \textbf{Falsifier:} the two curves are within each other's confidence
    intervals across the whole skew range, or FedAvg wins at high skew.
\end{tight}

\medskip
\textbf{Round 4 --- the one that makes it a contribution.} Round 3 may already
be known. A literature search (\chref{slr}) will tell you. Suppose it is known
for image data but not for tabular clinical data with missingness that is
itself hospital-dependent. \emph{That} is the delta, and now the claim is
yours:

\medskip
``On clinical tabular data where missingness is correlated with the
partitioning hospital, the FedProx advantage over FedAvg reported for image
benchmarks does not hold, because the proximal term penalises exactly the
local adaptation that hospital-specific missingness requires.''

\medskip
Note what happened: the claim acquired a \emph{mechanism}. It now predicts
something specific, which makes it far easier to refute --- and therefore far
more worth testing. All numbers here are illustrative; this is a shape to
imitate, not a result to cite.
\end{worked}

\armfig{
\begin{tikzpicture}[
  node distance=0pt,
  stage/.style={rounded corners=3pt, draw=primary!60, fill=primary!8,
                text width=2.55cm, align=center, font=\small,
                minimum height=1.25cm, inner sep=4pt},
  lbl/.style={font=\scriptsize\itshape, text=neutral, align=center,
              text width=2.5cm},
  ar/.style={-{Stealth[length=2.5mm]}, draw=secondary, thick}]

  \node[stage] (a) at (0,0) {\textbf{Subject area}\\[2pt]``federated learning''};
  \node[stage] (b) at (3.35,0) {\textbf{Question}\\[2pt]what is unknown here?};
  \node[stage] (c) at (6.7,0) {\textbf{Claim}\\[2pt]could be true or false};
  \node[stage, fill=secondary!12, draw=secondary!70]
        (d) at (10.05,0) {\textbf{Claim with mechanism}\\[2pt]predicts something};

  \draw[ar] (a) -- (b);
  \draw[ar] (b) -- (c);
  \draw[ar] (c) -- (d);

  \node[lbl, below=6pt of a] {no falsifier};
  \node[lbl, below=6pt of b] {a gap, found by\\reading};
  \node[lbl, below=6pt of c] {falsifiable, but\\maybe not new};
  \node[lbl, below=6pt of d] {falsifiable\\\emph{and} yours};

  \draw[ar, dashed, draw=accent] (d.north) .. controls +(0,1.0) and +(0,1.0) ..
        node[above, font=\scriptsize, text=accent]
        {evidence refutes it $\rightarrow$ revise, do not abandon} (b.north);
\end{tikzpicture}}
{From a subject area to a claim worth testing. The dashed return path is not
failure: a refuted claim that was sharp enough to refute has told you something,
and the revised question is better than the one you started with.}{claimpath}

\section{Novelty, and the Word ``Delta''}

Reviewers and examiners talk about the \term{delta}: the difference between
what was known before your work and what is known after. Three things are worth
knowing about it.

\textbf{The delta is measured against the literature, not against your own
prior ignorance.} Work is not novel because it is new to you. This is what
\chref{slr} is for, and it is why the review comes before the design.

\textbf{A small, well-established delta beats a large, unsupported one.} A
convincing demonstration that a widely believed result fails under one common
condition is a better contribution than an unvalidated framework claiming to
solve everything. Examiners have seen many of the latter.

\textbf{Negative and null results are deltas.} If a well-powered study finds no
effect where the field assumed one, that is knowledge. It is harder to publish,
which is a defect of the publication system (\chref{prereg}) rather than of the
result. Pre-registration is partly a device for protecting such findings.

\begin{regbox}[What your dissertation is judged on --- PhD Regulations 2024, \S14]
\S14(i): the dissertation ``must be an original and innovative contribution to
knowledge that contributes to solving socioeconomic problems''.

\medskip
\S14.1, on choosing the research area, requires that it:
\begin{tight}
  \item ``Correspond to the community needs at regional and local levels and
    comply with the priority national research agenda'';
  \item ``Reflect the basic and pure research'';
  \item ``Signify emerging areas of research that coincide preferably with
    sustainable development goals (SDGs)''.
\end{tight}
\end{regbox}

\begin{conceptbox}[Relevance is a constraint, not a decoration]
It is tempting to treat \S14.1 as a paragraph to add to the synopsis after the
real work is planned. Read it instead as a filter applied at the start: of the
several defensible claims you could test, prefer the one whose answer would
change a decision somebody in this region actually makes.

\medskip
That is also good research strategy, independent of the regulation. A claim
tied to a real setting comes with a real population to sample, real
practitioners to interview, and a real baseline to beat --- which is to say, it
comes with its evidence already half-specified. A claim floating free of any
setting has to invent all of that.
\end{conceptbox}

\section{Four Ways to Have No Contribution}

\begin{pitfall}[The four non-contributions]
\textbf{1. ``We applied X to Y.''} Applying an existing method to a new dataset
is not a contribution unless something about Y made the outcome uncertain. The
test: before running it, could you have confidently predicted the result? If
yes, you have done an exercise, not a study. The rescue is to find what makes Y
hard --- and that becomes the mechanism in your claim.

\medskip
\textbf{2. Benchmark chasing.} ``Our model reaches 94.2\%, beating the previous
94.0\%.'' Without an ablation showing \emph{which} change earned the 0.2, and
an interval showing it is not run-to-run noise, this is not a finding.
\chref{mleval} takes this apart.

\medskip
\textbf{3. The unevaluated tool.} A tool, a framework or an architecture,
presented with a feature list and screenshots. The claim is implicitly ``this
is good'', and nothing in the paper could have shown it was not.

\medskip
\textbf{4. The survey that is a list.} Sixty papers, tabulated by year and
venue, with a closing paragraph saying more research is needed. A review
contributes when it \emph{synthesises} --- reconciles disagreements, explains
why results differ, and identifies a gap precisely enough that somebody could
fill it.
\end{pitfall}

\begin{traditions}{a contribution}
\tradEMP{A contribution is a measured effect, reported with its uncertainty
and the conditions under which it was observed. The delta is a number
somebody did not have.}
\tradDSR{A contribution is an artefact plus the design knowledge it
demonstrates. The delta is a class of problem now known to be solvable this
way, and the principle that made it so.}
\tradQUAL{A contribution is an account that renders a setting intelligible ---
concepts, categories or a theory grounded in what participants actually do. The
delta is understanding, and it is judged by whether it illuminates, not by
whether it generalises.}
\tradFORM{A contribution is a proved property: a bound, an impossibility
result, a correctness argument, a complexity class. The delta is a theorem, and
the evidence is the proof.}
\end{traditions}

\begin{lab}[Classify five abstracts]
Take five recent papers from a venue in your area --- IEEE Xplore, the ACM
Digital Library or arXiv will do. For each, in no more than three lines:

\begin{tightnum}
  \item Which of the eight contribution types is it? (Some are two.)
  \item Write its claim as a single indicative sentence. If you cannot find one
    in the abstract, look in the introduction; if it is not there either, note
    that.
  \item What would have falsified it?
\end{tightnum}

\medskip
Then do the same for your own intended project. The discomfort you feel doing
step 3 for your own work is the point of the exercise. Keep this sheet --- it
becomes portfolio piece P1 (Appendix~\ref{app:schedule}).
\end{lab}

\begin{checkpoint}
\begin{tightnum}
  \item A colleague proposes: ``I will develop a machine learning model to
    predict student dropout at our university.'' Is this a contribution as
    stated? What single addition would make it one?
  \item Give an example of a result in computing that would be worth
    publishing even though the effect measured was zero.
  \item Why does a technique paper owe an ablation, while an empirical study
    owes a threats-to-validity section? What is each protecting the reader
    against?
  \item Your literature search shows your intended claim was published in 2019.
    Name three legitimate contributions still available to you.
\end{tightnum}
\end{checkpoint}

\begin{chaptersummary}
\begin{tight}
  \item A contribution is a claim, plus evidence that could have failed to
    support it, plus something the field did not already know.
  \item Development asks whether a thing can be built; research asks whether a
    claim is true. The same project can be either. An artefact contributes by
    carrying a claim, never by existing.
  \item Eight contribution types recur in computing, and each owes its reader a
    different kind of evidence. Most rejections are a mismatch between the type
    claimed and the evidence supplied.
  \item Write every project as claim, evidence, falsifier. An empty falsifier
    line means there is no project yet.
  \item The delta is measured against the literature, which is why the review
    comes first. A sharp small delta beats a vague large one.
  \item Your dissertation is additionally judged on relevance to community
    needs, the national research agenda and the SDGs. Treat that as a filter
    applied early, not a paragraph added late.
\end{tight}
\end{chaptersummary}

\begin{reviewq}
  \item Shaw argues that reviewers reject papers where question type, result
    type and validation type do not match. Construct a plausible computing
    paper in which each is individually reasonable but the three are mismatched,
    and say what a reviewer would object to.
  \item A design science researcher and an empiricist disagree about whether a
    novel architecture, evaluated only by a proof-of-concept implementation,
    is a contribution. Argue both positions, then state the condition under
    which you think each is right.
  \item ``Negative results are contributions'' is widely asserted and rarely
    acted on. Identify two structural reasons the publication system resists
    them, and one mechanism that partially corrects for this.
  \item \reg{PhD Regs 2024}{\S14.1} requires research areas to align with
    community needs and the national research agenda. Some would argue this
    constrains basic research, which the same clause also requires. Is there a
    genuine tension? Defend your answer with reference to a specific topic in
    your own area.
\end{reviewq}
